% APPENDIX D: Training Recipes and Failure Modes
% BASIS: Part II of the book -- Ch 2.6 ("Training Regimes, Optimization, and Failure
%   Modes"), whose numbers and mechanisms this file restates and never contradicts,
%   together with the three benchmark subsections of Ch 2.5 consolidated into it.
%   Regime names and grid sizes in the recipe table are
%   transcribed from Ch 2.5 through Appendix C's specification sheets, and the
%   weather row's K=2 is Ch 3.3's recorded two-step fine-tuning horizon.
% CONVENTION: "not recorded" marks silence in that record; nothing here
%   supplies a setting the record does not give.

\chapter{Training Recipes and Failure Modes}\label{app:d}

\chapref{ch:part2-training-regimes-optimization-and-failure-} argues; this appendix
lists. No number here is absent from that chapter, from the rest of Part~II, or
from the specification sheets of Appendix~\ref{app:c} that the first three rows
transcribe, save the weather row's $K=2$, named below. Where the record carries no
setting, the entry says so. Mechanisms are cited, not
re-derived (\eqref{eq:ch26-mode-gradient}, Table~\ref{tab:ch26-formats},
\eqref{eq:ch09-gronwall}). Nothing below is a measured outcome, and the qualification attached to any recipe
that is not a guarantee travels with it.
\index{training recipes!consolidated by regime}
\index{failure modes!quick reference}

% -----------------------------------------------------------------------------
\section{Schedule}
\label{sec:appd-schedule}
% -----------------------------------------------------------------------------

\chapref{ch:part2-training-regimes-optimization-and-failure-} states the symptom:
because complex multipliers modulate the modes directly, a constant learning rate
destabilizes early training. By
\eqref{eq:ch26-mode-gradient} gradient magnitude falls algebraically across one
array and is smallest at the band edge. The reference configuration,
Listing~\ref{lst:ch26-onecycle}, is \texttt{AdamW} at \texttt{weight\_decay=1e-4}
with \texttt{OneCycleLR} peaking at \texttt{max\_lr=5e-3}, rising from one
twenty-fifth of that, $2\times10^{-4}$, over the first fifth of the run
(\texttt{pct\_start=0.2}) and annealing to $2\times10^{-7}$.

Four settings carry it. \texttt{max\_lr} is the only absolute rate, so the
\texttt{lr=1e-3} in the \texttt{AdamW} call is a placeholder the first scheduler step
overwrites. \texttt{total\_steps} counts iterations and \texttt{scheduler.step()}
sits in the batch loop; mismatch them and the cycle ends early, the rate rests at
its floor, and the curve reads as a plateau. The \texttt{max\_norm=1.0} clip is a
global norm: one outlier gradient shrinks every tensor's update, and because Adam
divides out a clip that persists across steps, clipping guards against transients
rather than against an ill-scaled array. The annealing
law is cosine: band and schedule are one decision.

Two procedures follow. A sweep over mode budgets of eight, sixteen and twenty-four
must each hold \emph{steps per run proportional to the band}, since doubling it divides
the truncation contribution by roughly $2^{\,s-1/2}$ --- about $1.4$, $2.8$, $5.7$
at $s=1,2,3$ --- and if widening still moves nothing with steps scaled, the binding
term is capacity or the statistical and discretization contributions of
\cite{subedi2024controlling}. A port between grids requires rechecking rate and
decay for the reason of \chapref{ch:part2-training-regimes-optimization-and-failure-}:
the grid-dependent factor is a property of the stored coefficients, not of the
gradients the optimizer sees --- the inverse transform's adjoint cancels it exactly,
so switching the normalization convention changes what a reduced-precision format
must hold, not what \texttt{lr} or \texttt{weight\_decay} should be set to. The library workflow of
\chapref{ch:part2-fourier-neural-operators-fno-from-scratch} sets \texttt{lr=8e-3}
with no scheduler: an initial value, not a rate to be held constant.
\index{learning rate scheduling!reference configuration}
\index{gradient clipping!global norm}

% -----------------------------------------------------------------------------
\section{Numerical format}
\label{sec:appd-precision}
% -----------------------------------------------------------------------------

The symptom is a loss that becomes \texttt{NaN}, or gradients that become infinite,
a few hundred steps in on a graphics processing unit (GPU) with automatic mixed
precision (AMP) enabled, the same run being uneventful on the host central
processing unit (CPU).

The policy on record runs the $1\times1$ lifting and projection
convolutions in \texttt{bfloat16} and keeps the transform and the complex
multiplication in \texttt{float32}, and stops there.
\chapref{ch:part2-training-regimes-optimization-and-failure-} widens the
reduced-precision set to the bypass $W_\ell$ and the stored activations, on the
ground that nothing there accumulates. The protection must be an explicit cast
at the lift-to-spectral boundary: \texttt{torch.amp.autocast} casts the pointwise
convolutions and never the complex contraction, so the transform is exposed through
its input's dtype, not through an entry on a demotion list.

\texttt{float16} is excluded on range. Forward-unscaled, a unit-mean field on a
$256\times256$ grid has a direct-current coefficient of $65\,536$, thirty-two past
the largest finite binary16 value, and ensemble normalization does not repair it.
Sampling four times as finely in each direction multiplies a fixed-wavenumber
coefficient by about sixteen, while a steep spectrum takes modes near $k_{\max}$
below the least normal value $2^{-14}$. Where binary16 must be used, loss scaling
repairs the backward pass: multiply the loss by a large power of two, step, unscale.
\texttt{bfloat16} needs no scaler, its exponent matching \texttt{float32}.

Three cases are exclusions, not choices: a differenced residual, as in the
physics-informed neural operator (PINO) functional \eqref{eq:pino_loss}, where
cancellation and not range sets the derivative's error; the metric, since
\eqref{eq:rel_l2_loss} accumulated in \texttt{bfloat16} carries rounding no smaller
than its unit roundoff $2^{-8} = 3.9\times10^{-3}$, the order of the reported
floors; and binary16 gradient range. Losses and reported errors belong in single precision.
\index{automatic mixed precision (AMP)!policy recorded here}
\index{bfloat16!boundary cast at the spectral layer}

% -----------------------------------------------------------------------------
\section{Rollout horizon}
\label{sec:appd-rollout}
% -----------------------------------------------------------------------------

The trained object is the time-$\Delta t$ map \eqref{eq:ch26-rollout} applied to its
own output, and \eqref{eq:ch09-gronwall} amplifies one-step error at a rate set by
$\Lambda_\theta$: past one, the bound grows geometrically in the step count. The symptom is excellent single-step
accuracy with a rollout that dies around step twenty.

The recipe is the pushforward objective \eqref{eq:ch26-pushforward}, paid for by
truncated backpropagation --- unroll $K$ steps, detach the intermediate states,
keep gradients on the final segment only --- since one step of the full sum costs
$\mathcal{O}(K)$ transforms and, if differentiated exactly, $\mathcal{O}(K)$ stored
activations. The mechanism is a change of
sampling distribution, not a new penalty: the terms do not bound $\Lambda_\theta$, no
theorem here relates $\lambda_k$ or $K$ to \eqref{eq:ch09-lipschitz}, and detaching
leaves composed stability only weakly supervised, the network accurate \emph{at} the
states it reaches and not \emph{by the route}.

The inherited record states that $\lambda_k$ must decay geometrically or $K$ be
capped near the horizon the physics allows, uniform $\lambda_k$ being the aggressive
setting; it gives no value of any $\lambda_k$, and the horizon the remedy names is the
word ``modest''; where the book attaches a number to an unrolled objective, the
weather row below carries it --- $K=2$, a two-step fine-tuning. The check precedes
them: evaluate
\eqref{eq:ch09-lipschitz} at the checkpoint, which Proposition~\ref{prop:ch09-lip}
bounds from the parameters; if it exceeds one and grows, the objective is satisfied
and the dynamics are not. Then, in increasing cost, reduce the peak rate or deepen
the anneal, add pushforward terms at modest $K$, and ask whether the model collapsed
to a smooth state rather than diverged --- the conditional-mean failure, judged by
the weak error \eqref{eq:ch09-weak}. Pushforward's inputs are finitely many
trajectories, so a drifting rollout outruns its frozen normalizer.
\index{pushforward regularization!horizon and weights as recorded}
\index{truncated backpropagation!detached rollout states}

% -----------------------------------------------------------------------------
\section{Two checks and one diagnosis}
\label{sec:appd-checks}
% -----------------------------------------------------------------------------

Bookkeeping explains a floor that moves when only the batch size or device count
changed. The per-sample normalization of \eqref{eq:rel_l2_loss} weights
low-amplitude members heavily. Then \texttt{total\_steps} must equal the calls
\texttt{scheduler.step()} will actually make: recompute it and the cosine keeps its
shape while the update count halves, leave it stale and the run stops mid-descent
with the rate still high. Only the second is a mistake, caught by logging the rate
and not the loss. A batch study must fix the step count or it measures the
schedule, not the batch.

The second check surveils the spectral branch, the failure that leaves every curve
falling. Plot $\sup_k \lVert R_{\phi,\ell}(k)\rVert$ against $\lVert W_\ell\rVert$
per layer per epoch and reduce the decay coefficient if the ratio trends downward:
decay shrinks a multiplier's magnitude and leaves its phase untouched, so a
driven-small branch leaves the bypass as the layer's only content, a shallow
feedforward map without the global coupling that made it an operator.

The last is a diagnosis: a plausible but visibly smoothed output on a refined grid
is not a broken model, since refinement raises the truth's Nyquist frequency while
the band addresses only the fixed wavenumbers of \eqref{eq:ch05-modeset}. Reported
super-resolution assumes inputs band-limited inside the retained set, so it demonstrates resolution transfer and not
recovery of lost detail, and \eqref{eq:limit_commute} should be read with that premise
attached; if content above $k_{\max}$ is needed, widen the band and retrain.
\index{weight decay!spectral-branch surveillance}
\index{resolution transfer!diagnosis of a smoothed output}

% -----------------------------------------------------------------------------
\section{Recipes by regime}
\label{sec:appd-table}
% -----------------------------------------------------------------------------

Table~\ref{tab:appd-recipes} collects these entries by regime, naming the first
three rows after the specifications of Appendix~\ref{app:c} and the last after a
port rather than a problem. Take its conclusion before the cells: the
schedule column repeats because Part~II supplies one configuration and not six, and
the regimes separate on three axes --- whether the map is consumed once or composed,
whether the transform sits inside the reduced-precision region, and whether a
differenced residual enters the objective at all.

\begin{table}[t]
\centering
\footnotesize
\caption{Training recipes consolidated by regime. Every cell is stated in
\chapref{ch:part2-training-regimes-optimization-and-failure-}, elsewhere in Part~II, or in
the specification sheets of Appendix~\ref{app:c} that the first three rows transcribe,
and ``not recorded'' marks silence there; the one exception, $K=2$ on the weather
row, is the fine-tuning horizon of
\chapref{ch:part3-global-weather-forecasting-and-climate-dyn}. The schedule column
repeats because only one reference configuration is on record
(Listing~\ref{lst:ch26-onecycle}), not a schedule per benchmark.}
\label{tab:appd-recipes}
\begin{tabularx}{\textwidth}{@{}>{\raggedright\arraybackslash}p{0.165\textwidth}
  >{\raggedright\arraybackslash}p{0.125\textwidth}
  >{\raggedright\arraybackslash}p{0.195\textwidth}
  >{\raggedright\arraybackslash}p{0.18\textwidth}
  >{\raggedright\arraybackslash}X@{}}
\toprule
\textbf{Regime} & \textbf{Schedule} & \textbf{Precision} & \textbf{Rollout strategy} & \textbf{Known failure and first check} \\
\midrule
Steady coefficient-to-state, $64^2$ (\eqref{eq:ch25-darcy})
& Listing~\ref{lst:ch26-onecycle}
& $1\times1$ maps \texttt{bfloat16}; transform and product \texttt{float32}
& None: no time step
& Flat across mode budgets: scale steps with the band before blaming truncation \\
\addlinespace
Evolution map on $t\in[0,1]$, $N=128$ (\eqref{eq:ch25-burgers})
& Same; a constant rate destabilizes early
& Same policy; one-dimensional transform
& Not recorded: one model per $t$, $t$ as a channel, or an iterated map
& Early instability on a fixed rate \\
\addlinespace
Time-marching map rolled fifty to one hundred steps (\eqref{eq:ch25-navierstokes})
& Same; no per-benchmark rate recorded
& Same policy; on \texttt{float16}, add loss scaling
& Pushforward \eqref{eq:ch26-pushforward}, truncated backpropagation, detached states, geometric $\lambda_k$, ``modest'' $K$
& Dies near step twenty: compute \eqref{eq:ch09-lipschitz} at the checkpoint first \\
\addlinespace
Planetary-scale rollout, six-hour steps over days \cite{pathak2022fourcastnet}
& Not recorded
& Not recorded
& Same objective at $K=2$; a horizon past predictability is the failure
& Bounded, smooth and wrong: the conditional mean, judged by \eqref{eq:ch09-weak} \\
\addlinespace
Physics-informed residual objective (\eqref{eq:pino_loss})
& Not recorded
& Residual and metric outside reduced precision
& Not recorded
& Derivative error set by the format; a wider exponent does not help \\
\addlinespace
Any of the above, ported to a refined grid
& Recheck peak rate and decay
& \texttt{float16} overflows: direct-current term past $65\,504$ on $256^2$
& Unchanged; no refinement-specific setting recorded
& \texttt{NaN} on GPU with AMP, clean on CPU: act on the range, not the model \\
\bottomrule
\end{tabularx}
\end{table}
